\section{Top-2 NFT Migration and the Original NFT Liquidity Protocol}
\label{sec:nft-migration}

The \Zoo{} \DEX{} is the migration destination for \NFT{} markets,
and the realisation of the original ``NFT Liquidity Protocol''
coinage from the October 31, 2021 \Zoo{} whitepaper \cite{zoo-2021}.

\subsection{What the 2021 paper coined}

The 2021 \Zoo{} whitepaper introduced ``NFT Liquidity Protocol'' as
the term for a system in which \NFT{}s become productive collateral
inside a yield-bearing on-chain market. The 2021 vision had three
moving parts:

\begin{enumerate}[leftmargin=1.4em]
  \item Wildlife conservation \NFT{}s as on-chain collateral.
  \item Borrowing and lending against the \NFT{}s without the
        holder having to sell.
  \item Yield mechanisms that fed back to conservation.
\end{enumerate}

The original BSC implementation in 2021 had infrastructural
limitations (no native fractional secondary trading, no on-chain
appraisal, no privacy primitives, no triple-cert finality). The
\Zoo{} \DEX{} on \Zoo{} D-Chain in April 2026 is the first venue
where every original 2021 commitment is implementable end-to-end.

\subsection{Top-2 NFT migration plan}

\Zoo{} 4.0 launched (February 14, 2026) with an explicit go-to-market
for \NFT{} migration, targeting the leading \NFT{} collections in
wildlife/conservation, AI/agent, and metaverse categories
\cite{zoo-4-0}. The \Zoo{} \DEX{} launch on April 20, 2026 makes that
go-to-market actionable: the migrating collections now have a venue
on which to trade and to borrow against.

\subsection{Native fractional NFT trading}

A whole \NFT{} on \Zoo{} D-Chain can be wrapped into a Fractional
ERC-4626 vault that issues fungible shares. The vault shares trade
on \Zoo{} \DEX{} V3 as a concentrated-liquidity pair against
\$ZOO or stablecoin. Properties:

\begin{itemize}[leftmargin=1.2em]
  \item \textbf{Continuous price discovery.} The vault share's price
        is a continuous market quote, not an episodic auction.
  \item \textbf{Composable.} The vault shares are ZRC-20 tokens;
        they compose into AMM pools, lending markets, derivatives,
        and the \Liquidity{}-precompile-protected transfer flow.
  \item \textbf{Re-aggregable.} Holders accumulating sufficient
        shares can redeem the underlying \NFT{} from the vault by
        burning all outstanding shares.
\end{itemize}

\subsection{Collateral-backed lending (the realised 2021 vision)}

A holder of a high-value \NFT{} can use it as collateral for a loan
in stablecoin, \$ZOO, or any listed asset. The flow:

\begin{enumerate}[leftmargin=1.4em]
  \item Holder locks the \NFT{} in a loan-escrow contract.
  \item Loan terms (principal, APR, duration, collateral floor)
        are agreed with a lender (or pooled-lender market).
  \item Counterparties sign the loan agreement under the post-quantum
        scheme inherited from \Zoo{} 3.0 \cite{zoo-3-0} (\Corona{}
        + $\MLDSA{}\text{-}65$).
  \item Holder receives the loan principal in the chosen denomination.
  \item Liquidation oracle (A-Chain attested) monitors the
        collateral floor against the appraisal feed (which can run
        confidentially via F-Chain \cite{lp-013}).
  \item Holder repays principal plus interest by the maturity date;
        the \NFT{} is released. Or the collateral floor is breached
        and liquidation triggers per \Zoo{} 3.0 NFT Liquidity
        Protocol semantics.
\end{enumerate}

This is exactly the 2021 paper's vision, with the missing
infrastructural pieces filled in by \Zoo{} 3.0 (PQ signing) and
\Zoo{} 4.0 (\GPU{}-native execution and \DEX{}).

\subsection{Provenance for AI/agent NFT collections}

Collections whose value derives from on-chain AI provenance
(generative-art collections, agent \NFT{}s, on-chain creatures with
AI behaviour) benefit from A-Chain-anchored TEE inference receipts
\cite{lp-134}. The provenance chain on \Zoo{} D-Chain is verifiable
end-to-end: every act of generation that produced a token in the
collection is traceable through the inference-receipt chain to a
TEE-attested execution.

\subsection{Token-bound accounts}

ZRC-721 \NFT{}s on \Zoo{} D-Chain have native ERC-6551-style
token-bound accounts. The account can hold \$ZOO, per-\LLM{} tokens,
other \NFT{}s, and accrue yield from staking or LP positions.
Combined with the \Zoo{} \DEX{}, the \NFT{} becomes a programmable
agent: it can hold a portfolio, execute trades on its own behalf
within whitelisted bounds, and participate in governance.

\subsection{Migration mechanics}

\begin{enumerate}[leftmargin=1.4em]
  \item \textbf{Snapshot.} \Zoo{} Bridge takes a snapshot of holders
        on the source chain (Ethereum, BSC, etc.).
  \item \textbf{Lock-and-mint.} The source-chain contract is escrowed
        via the bridge committee; the \Zoo{} D-Chain native ZRC-721
        is minted to the same wallet address.
  \item \textbf{Metadata mirror.} On-chain metadata moves over with
        provenance hash chain intact; off-chain assets re-pin to
        \Zoo{}'s decentralised storage layer.
  \item \textbf{Treasury onboarding.} The collection's treasury is
        onboarded to a \Zoo{} \DAO{}-governed multisig with optional
        F-Chain confidential balance.
  \item \textbf{Royalty redirection.} Future royalties accrue to the
        D-Chain contract; the bridge keeps the source-chain contract
        readable for legacy explorers.
\end{enumerate}

The bridge committee under \Zoo{} 3.0 PQ discipline \cite{zoo-3-0}
attests the migration; the destination mint carries triple-cert
finality.

\subsection{Continuous secondary trading}

After migration, every transfer of an \NFT{} from the migrated
collection trades on \Zoo{} \DEX{}: as a whole-token transfer
against a price-discovery pair, as a fractional share via the V3
vault pool, or as collateral inside a loan agreement. The
collection's secondary-market liquidity is materially better on
\Zoo{} \DEX{} than on its source chain because of:

\begin{itemize}[leftmargin=1.2em]
  \item Sub-1.1s settlement (vs. multi-block confirmation
        elsewhere).
  \item Native fractional secondary market.
  \item Composable loan agreements.
  \item Confidential balance and orderflow optionality.
\end{itemize}

\subsection{Why this closes a 2021 commitment}

The 2021 \Zoo{} whitepaper coined ``NFT Liquidity Protocol'' but
the BSC infrastructure of the time could not deliver the full
promise. \Zoo{} \DEX{} on \Zoo{} D-Chain in April 2026 delivers it:
sub-second-final, post-quantum, compliance-perimetered, with
native fractional secondary, confidential balances, and a
collateralised lending primitive. This paper is the activation
document for that closure.
